%% ----------------------------------------------------------------------------
%%  §7.5  Massimi e minimi relativi
%% ----------------------------------------------------------------------------
\section{Massimi e minimi relativi}
\label{sec:estremi-due}

Le definizioni di massimo e di minimo relativo si estendono senza difficoltà,
sostituendo gli intervalli con gli intorni circolari.

\begin{definizione}[Massimi e minimi relativi]\label{def:estremi-due}
\index{massimo!relativo!in più variabili}\index{minimo!relativo!in più variabili}
Sia $f$ una funzione di due variabili definita in un insieme $D\subseteq\R^{2}$. Un punto
$(x_{0},y_{0})\in D$ è di \emph{massimo relativo} per $f$ in $D$ se esiste un intorno
circolare $I_{\delta}$ di $(x_{0},y_{0})$ tale che
\begin{equation}\label{eq:max-due}
  f(x_{0},y_{0})\ge f(x,y) \qquad\perogni(x,y)\in I_{\delta}\cap D ;
\end{equation}
è di \emph{minimo relativo} se $f(x_{0},y_{0})\le f(x,y)$ per ogni
$(x,y)\in I_{\delta}\cap D$. Se le disuguaglianze valgono in senso stretto per
$(x,y)\neq(x_{0},y_{0})$, il massimo o il minimo si dice \emph{stretto}.
\end{definizione}

Come per le funzioni di una variabile, la derivata dà informazioni soltanto nei punti
«interni».

\begin{definizione}[Punto interno]\label{def:punto-interno}\index{punto!interno}
Un punto $(x_{0},y_{0})$ è \emph{interno} all'insieme $D\subseteq\R^{2}$ se esiste un
intorno circolare $I_{\delta}$ di $(x_{0},y_{0})$ contenuto in $D$. Un insieme i cui punti
sono tutti interni si dice \emph{aperto}.
\end{definizione}

La \cref{fig:punto-interno} illustra la differenza tra i punti interni e quelli che
stanno sul bordo dell'insieme.

\begin{figure}[htb]
\centering
\begin{tikzpicture}[scale=1.0]
  \draw[-Latex, grigio] (-0.3,0) -- (6.4,0) node[right, black] {$x$};
  \draw[-Latex, grigio] (0,-0.3) -- (0,3.6) node[above, black] {$y$};
  \fill[pattern=north east lines, pattern color=blu]
    (0.8,3.2) .. controls (0.6,1.2) and (2.0,0.1) .. (3.5,0.6)
    .. controls (5.0,1.1) and (6.0,2.4) .. (5.9,3.3) -- (0.8,3.2) -- cycle;
  \draw[blu, thick]
    (0.8,3.2) .. controls (0.6,1.2) and (2.0,0.1) .. (3.5,0.6)
    .. controls (5.0,1.1) and (6.0,2.4) .. (5.9,3.3);
  \fill[white] (2.2,2.0) circle (0.75);
  \fill[pattern=north east lines, pattern color=blu!50] (2.2,2.0) circle (0.75);
  \draw[rame, thick, dashed] (2.2,2.0) circle (0.75);
  \fill (2.2,2.0) circle (1.5pt) node[above right, font=\scriptsize] {$(x_{0},y_{0})$};
  \node[font=\scriptsize, fill=white, inner sep=1pt] at (2.2,1.6) {punto interno};
  \fill (5.0216,1.507) circle (1.5pt); % punto della curva di Bézier (t = 0,4)
  \node[font=\scriptsize, right] at (5.25,1.25) {punto non interno};
  \node[font=\scriptsize, blu] at (4.6,2.6) {$D$};
\end{tikzpicture}
\caption{Un punto interno all'insieme $D$, che ha un intorno circolare tutto contenuto in
$D$, e un punto di $D$ che sta sul bordo, ogni intorno del quale contiene punti esterni a
$D$.}
\label{fig:punto-interno}
\end{figure}

\begin{teorema}[Condizione necessaria per gli estremi]\label{teo:fermat-due}
\index{teorema!di Fermat!in più variabili}
Se $(x_{0},y_{0})$ è un punto di massimo o di minimo relativo interno al dominio $D$ di una
funzione $f(x,y)$, e se $f$ è dotata di derivate parziali prime in $(x_{0},y_{0})$, allora
\begin{equation}\label{eq:fermat-due}
  f_{x}(x_{0},y_{0})=0, \qquad f_{y}(x_{0},y_{0})=0 .
\end{equation}
\end{teorema}

\begin{schema}
\Ipotesi{$(x_{0},y_{0})$ interno a $D$, di massimo (o minimo) relativo; esistono
$f_{x}(x_{0},y_{0})$, $f_{y}(x_{0},y_{0})$.}
\Tesi{$\nabla f(x_{0},y_{0})=(0,0)$.}
\end{schema}

\begin{strategia}
Restringiamo $f$ alle rette per $(x_{0},y_{0})$ parallele agli assi. Si ottengono funzioni
di una variabile che hanno un estremo relativo interno, e alle quali si applica il
teorema di Fermat; le loro derivate sono proprio le derivate parziali.
\end{strategia}

\begin{dimostrazione}
Supponiamo che $(x_{0},y_{0})$ sia un punto di massimo relativo interno; esiste allora un
intorno circolare $I_{\delta}$, contenuto in $D$, nel quale vale la \eqref{eq:max-due}.
Fissato $y=y_{0}$, consideriamo la funzione di una variabile reale
$F(x)=f(x,y_{0})$, definita nell'intervallo $(x_{0}-\delta,x_{0}+\delta)$ (i punti
$(x,y_{0})$ con $\abs{x-x_{0}}<\delta$ stanno in $I_{\delta}$). Per la \eqref{eq:max-due},
$F(x_{0})\ge F(x)$ per ogni $x$ di tale intervallo: $x_{0}$ è un punto di massimo
relativo interno per $F$. Per il teorema di Fermat (\cref{teo:fermat}),
$F'(x_{0})=f_{x}(x_{0},y_{0})=0$. Si procede in modo analogo, con la funzione
$G(y)=f(x_{0},y)$, per ottenere $f_{y}(x_{0},y_{0})=0$.
\end{dimostrazione}

I punti di massimo o di minimo interni al dominio vanno dunque ricercati tra i punti che
annullano entrambe le derivate parziali prime; tali punti si dicono \emph{punti critici}
(o \emph{stazionari}) della funzione: sono i punti in cui il gradiente si annulla. Come
per le funzioni di una variabile, la condizione è necessaria ma non sufficiente: la
funzione $y^{2}-x^{2}$ ha gradiente $(-2x,2y)$, nullo nell'origine, ma l'origine non è né
di massimo né di minimo, perché $f(x,0)=-x^{2}<0$ e $f(0,y)=y^{2}>0$ per $x,y\neq0$. Un
punto critico di questo tipo si chiama \emph{punto di sella}, per la forma del grafico
(\cref{fig:superfici}).

\subsection{Il criterio dell'hessiano}

Il teorema che segue fornisce un criterio per stabilire la natura di un punto critico.
Supponiamo che $f$ abbia derivate seconde continue in un insieme aperto $D$; si definisce
il \emph{determinante hessiano} (o semplicemente \emph{hessiano})
\begin{equation}\label{eq:hessiano}
  H(x,y)=\begin{vmatrix} f_{xx}(x,y) & f_{xy}(x,y)\\ f_{yx}(x,y) & f_{yy}(x,y)\end{vmatrix}
  =f_{xx}(x,y)\,f_{yy}(x,y)-\bigl(f_{xy}(x,y)\bigr)^{2},
\end{equation}
dove si è tenuto conto che, per il teorema di Schwarz, $f_{xy}=f_{yx}$. Talvolta, per
mettere in luce la dipendenza dalla funzione $f$, si scrive $Hf(x,y)$.

\begin{teorema}[Criterio dell'hessiano]\label{teo:hessiano}\index{criterio!dell'hessiano}
\index{hessiano}
Sia $f$ dotata di derivate parziali prime e seconde continue in un intorno del punto
$(x_{0},y_{0})$, e sia $(x_{0},y_{0})$ un punto critico:
$f_{x}(x_{0},y_{0})=f_{y}(x_{0},y_{0})=0$.
\begin{enumerate}[label=\textup{(\roman*)}]
  \item Se $Hf(x_{0},y_{0})>0$ e $f_{xx}(x_{0},y_{0})>0$, allora $(x_{0},y_{0})$ è un punto di
  minimo relativo (stretto).
  \item Se $Hf(x_{0},y_{0})>0$ e $f_{xx}(x_{0},y_{0})<0$, allora $(x_{0},y_{0})$ è un punto di
  massimo relativo (stretto).
  \item Se $Hf(x_{0},y_{0})<0$, allora $(x_{0},y_{0})$ non è né di massimo né di minimo
  relativo: si dice che è un \emph{punto di sella}.
\end{enumerate}
\end{teorema}

\begin{schema}
\Ipotesi{$f$ con derivate prime e seconde continue vicino a $(x_{0},y_{0})$;
$\nabla f(x_{0},y_{0})=(0,0)$.}
\Tesi{$H>0$, $f_{xx}>0$: minimo; \ $H>0$, $f_{xx}<0$: massimo; \ $H<0$: sella (tutto
calcolato in $(x_{0},y_{0})$).}
\end{schema}

\begin{strategia}
Restringiamo $f$ al segmento che congiunge $(x_{0},y_{0})$ con un punto vicino
$(x_{0}+h,y_{0}+k)$: si ottiene una funzione di una variabile $g(t)$ con $g'(0)=0$, la cui
derivata seconda è la \emph{forma quadratica}
$f_{xx}h^{2}+2f_{xy}hk+f_{yy}k^{2}$. Il segno di questa forma, per tutte le direzioni $(h,k)$,
è deciso proprio da $f_{xx}$ e da $H$ (completando il quadrato). Per rendere rigoroso il
ragionamento occorrono due strumenti: la derivata di $f$ lungo un segmento, che richiede
la nozione di differenziabilità, e una formula di Taylor con il resto nella forma di
Lagrange. Per questo la dimostrazione è rinviata al \cref{sec:dim-hessiano}, in
appendice.
\end{strategia}

Un punto critico è dunque di massimo o di minimo relativo se $Hf(x_{0},y_{0})>0$, mentre
se $Hf(x_{0},y_{0})<0$ non è né di massimo né di minimo. Il caso $Hf(x_{0},y_{0})=0$ non è
contemplato nell'enunciato, perché in tal caso il punto può essere di massimo, di minimo
o di sella: nell'origine le funzioni $x^{4}+y^{4}$, $-(x^{4}+y^{4})$ e
$x^{4}-y^{4}$ hanno tutte $H=0$, ma la prima ha un minimo, la seconda un massimo, la terza
una sella. Si noti infine che, se $H>0$, allora $f_{xx}f_{yy}>f_{xy}^{2}\ge0$, e quindi
$f_{xx}$ e $f_{yy}$ hanno lo stesso segno: nel criterio si può usare l'una o l'altra.

\begin{esempio}[Quattro funzioni e i loro punti critici]\label{es:punti-critici}
\begin{enumerate}[label=\textup{(\alph*)}]
  \item La funzione $f(x,y)=x^{2}+y^{2}$ ha derivate parziali $f_{x}=2x$, $f_{y}=2y$, che si
  annullano nel punto $(0,0)$. L'hessiano vale
  $H=\begin{vmatrix}2&0\\0&2\end{vmatrix}=4>0$ ed $f_{xx}=2>0$: l'origine è di minimo
  relativo. Si verifica anche direttamente che è di minimo \emph{assoluto}, perché
  $f(0,0)=0\le x^{2}+y^{2}$ per ogni $(x,y)$.
  \item La funzione $f(x,y)=y^{2}-x^{2}$ ha l'unico punto critico $(0,0)$, con
  $H=\begin{vmatrix}-2&0\\0&2\end{vmatrix}=-4<0$: è un punto di sella, ben riconoscibile
  nel grafico del paraboloide iperbolico.
  \item La funzione $f(x,y)=x^{3}-6y(x+y)=x^{3}-6xy-6y^{2}$, definita su $\R^{2}$, ha i
  punti critici che risolvono il sistema
  \[
    f_{x}=3x^{2}-6y=0, \qquad f_{y}=-6x-12y=0 .
  \]
  Dalla seconda equazione $x=-2y$, che sostituito nella prima dà $12y^{2}-6y=0$, cioè
  $y=0$ oppure $y=\frac12$; in corrispondenza $x=0$ oppure $x=-1$. I punti critici sono
  $(0,0)$ e $(-1,\frac12)$. Poiché $f_{xx}=6x$, $f_{xy}=-6$, $f_{yy}=-12$,
  \[
    H(x,y)=\begin{vmatrix}6x&-6\\-6&-12\end{vmatrix}=-72x-36 .
  \]
  Nell'origine $H(0,0)=-36<0$: è un punto di sella. Nel punto $(-1,\frac12)$ si ha
  $H=36>0$ ed $f_{xx}(-1,\frac12)=-6<0$: è un punto di massimo relativo, con
  $f(-1,\frac12)=\frac12$ (\cref{fig:punti-critici}, a sinistra). Non è un massimo
  assoluto, perché $f(x,0)=x^{3}\to+\infty$ per $x\to+\infty$.
  \item Per $f(x,y)=y\,e^{-(x^{2}+y^{2})}$ il sistema
  \[
    f_{x}=-2xy\,e^{-(x^{2}+y^{2})}=0, \qquad f_{y}=(1-2y^{2})\,e^{-(x^{2}+y^{2})}=0
  \]
  dà, dalla prima equazione, $x=0$ oppure $y=0$; ma $y=0$ è incompatibile con la seconda,
  che richiede $y=\pm\frac{\sqrt2}2$. I punti critici sono $(0,\pm\frac{\sqrt2}2)$. Il
  programma seguente calcola l'hessiano: in entrambi i punti $H=4e^{-1}>0$, mentre
  $f_{xx}(0,\frac{\sqrt2}2)=-\sqrt2\,e^{-1/2}<0$ e $f_{xx}(0,-\frac{\sqrt2}2)=\sqrt2\,e^{-1/2}>0$:
  il primo è di massimo, il secondo di minimo (\cref{fig:punti-critici}, a destra).
\end{enumerate}
\end{esempio}

\begin{figure}[htb]
\centering
\includegraphics[width=\textwidth]{punti_critici}
\caption{A sinistra, le curve di livello di $x^{3}-6y(x+y)$: intorno al massimo
$(-1,\frac12)$ sono curve chiuse; nella sella $(0,0)$ si incrociano due rami della curva
di livello $0$ (in grigio). A destra, il grafico di $y\,e^{-(x^{2}+y^{2})}$ con il massimo e
il minimo (\texttt{python/fig\_punti\_critici.py}).}
\label{fig:punti-critici}
\end{figure}

\begin{python}[Ricerca e classificazione dei punti critici]
import sympy as sp

x, y = sp.symbols("x y", real=True)

def classifica(f):
    """Punti critici di f(x, y) e loro natura con l'hessiano."""
    fx, fy = sp.diff(f, x), sp.diff(f, y)
    H = sp.hessian(f, (x, y))
    for P in sp.solve([fx, fy], [x, y], dict=True):
        h = sp.simplify(H.det().subs(P))
        fxx = sp.simplify(H[0, 0].subs(P))
        if h > 0:
            tipo = "minimo" if fxx > 0 else "massimo"
        elif h < 0:
            tipo = "sella"
        else:
            tipo = "H = 0: il criterio non decide"
        print(f"({P[x]}, {P[y]}): H = {h}, f_xx = {fxx} -> {tipo}")

for f in (x**2 + y**2, y**2 - x**2, x**3 - 6*y*(x + y),
          y * sp.exp(-(x**2 + y**2)), x**3 + y**3 - 3*x*y):
    print("f =", f)
    classifica(f)
\end{python}
\begin{risultato}
f = x**2 + y**2
(0, 0): H = 4, f_xx = 2 -> minimo
f = -x**2 + y**2
(0, 0): H = -4, f_xx = -2 -> sella
f = x**3 - 6*y*(x + y)
(-1, 1/2): H = 36, f_xx = -6 -> massimo
(0, 0): H = -36, f_xx = 0 -> sella
f = y*exp(-x**2 - y**2)
(0, -sqrt(2)/2): H = 4*exp(-1), f_xx = sqrt(2)*exp(-1/2) -> minimo
(0, sqrt(2)/2): H = 4*exp(-1), f_xx = -sqrt(2)*exp(-1/2) -> massimo
f = x**3 - 3*x*y + y**3
(0, 0): H = -9, f_xx = 0 -> sella
(1, 1): H = 27, f_xx = 6 -> minimo
\end{risultato}

L'ultima funzione dell'elenco, $x^{3}+y^{3}-3xy$, è oggetto di un esercizio del
\cref{sec:esercizi-cap7}. Chiudiamo gli esempi con un punto di massimo in cui
l'hessiano è nullo. La funzione $f(x,y)=\cos(x^{2}+y^{2})$ della \cref{fig:superfici} ha
nell'origine un massimo assoluto, perché $f(0,0)=\cos0=1\ge\cos(x^{2}+y^{2})$ per ogni
$(x,y)$; in accordo con il \cref{teo:fermat-due}, l'origine annulla il gradiente
$\bigl(-2x\sen(x^{2}+y^{2}),\,-2y\sen(x^{2}+y^{2})\bigr)$. Ma tutte le derivate seconde,
\begin{align*}
  f_{xx}&=-2\sen(x^{2}+y^{2})-4x^{2}\cos(x^{2}+y^{2}),\\
  f_{yy}&=-2\sen(x^{2}+y^{2})-4y^{2}\cos(x^{2}+y^{2}),\\
  f_{xy}&=-4xy\cos(x^{2}+y^{2}),
\end{align*}
si annullano nell'origine, e così l'hessiano: il criterio non è applicabile, e la natura
del punto si stabilisce con il confronto diretto.

\subsection{Un'applicazione: il metodo dei minimi quadrati}

In galleria del vento si misura il coefficiente di resistenza $C_{D,i}$ di un modello per
alcuni valori $C_{L,i}$ del coefficiente di portanza, $i=1,\dots,n$. Vogliamo stimare i
parametri della polare parabolica $C_{D}=C_{D0}+kC_{L}^{2}$. Posto $X_{i}=C_{L,i}^{2}$, si
tratta di trovare la retta $C_{D}=b+aX$ che meglio approssima i punti $(X_{i},C_{D,i})$.
Il criterio più usato, dovuto a Legendre e a Gauss, è quello dei \emph{minimi
quadrati}: si scelgono $a$ e $b$ in modo da rendere minima la somma dei quadrati degli
scarti (o \emph{residui})
\begin{equation}\label{eq:minimi-quadrati}
  S(a,b)=\sum_{i=1}^{n}\bigl(C_{D,i}-b-aX_{i}\bigr)^{2} .
\end{equation}
È una funzione di due variabili su tutto $\R^{2}$. I suoi punti critici risolvono il
sistema
\[
  S_{a}=-2\sum_{i=1}^{n}X_{i}\bigl(C_{D,i}-b-aX_{i}\bigr)=0, \qquad
  S_{b}=-2\sum_{i=1}^{n}\bigl(C_{D,i}-b-aX_{i}\bigr)=0,
\]
cioè le \emph{equazioni normali}
\begin{equation}\label{eq:equazioni-normali}
  a\sum_{i}X_{i}^{2}+b\sum_{i}X_{i}=\sum_{i}X_{i}C_{D,i}, \qquad
  a\sum_{i}X_{i}+b\,n=\sum_{i}C_{D,i} .
\end{equation}
È un sistema lineare il cui determinante è $n\sum_{i}X_{i}^{2}-\bigl(\sum_{i}X_{i}\bigr)^{2}$.
Questa quantità è positiva se i valori $X_{i}$ non sono tutti uguali, grazie all'identità
\begin{equation}\label{eq:identita-lagrange}
  n\sum_{i}X_{i}^{2}-\Bigl(\sum_{i}X_{i}\Bigr)^{2}=\frac12\sum_{i=1}^{n}\sum_{j=1}^{n}(X_{i}-X_{j})^{2},
\end{equation}
che si verifica sviluppando il quadrato a secondo membro. Il sistema ha quindi
un'unica soluzione $(a^{*},b^{*})$. Le derivate seconde sono costanti:
$S_{aa}=2\sum_{i}X_{i}^{2}>0$, $S_{ab}=2\sum_{i}X_{i}$, $S_{bb}=2n$, e l'hessiano
$H=4\bigl[n\sum_{i}X_{i}^{2}-(\sum_{i}X_{i})^{2}\bigr]$ è positivo: per il criterio
dell'hessiano $(a^{*},b^{*})$ è un punto di minimo relativo.

In questo caso si può dire di più: il minimo è \emph{assoluto}. Indichiamo con
$r_{i}=C_{D,i}-b^{*}-a^{*}X_{i}$ i residui nel punto di minimo; le equazioni normali dicono
che $\sum_{i}r_{i}=0$ e $\sum_{i}X_{i}r_{i}=0$. Allora, per ogni $(\alpha,\beta)$,
\begin{align*}
  S(a^{*}+\alpha,b^{*}+\beta)&=\sum_{i}\bigl(r_{i}-\alpha X_{i}-\beta\bigr)^{2}\\
  &=S(a^{*},b^{*})-2\alpha\sum_{i}X_{i}r_{i}-2\beta\sum_{i}r_{i}+\sum_{i}(\alpha X_{i}+\beta)^{2},
\end{align*}
e i due termini centrali sono nulli: $S(a^{*}+\alpha,b^{*}+\beta)\ge S(a^{*},b^{*})$, con
l'uguaglianza soltanto per $\alpha=\beta=0$. Il programma applica il metodo a un insieme
di misure simulate.

\begin{python}[Stima di $C_{D0}$ e $k$ con i minimi quadrati]
import numpy as np

# "prove in galleria": C_D misurato per vari C_L (dati simulati)
rng = np.random.default_rng(7)
CL = np.linspace(0.0, 1.2, 13)
CD = 0.025 + 0.0497 * CL**2 + rng.normal(0, 0.0008, CL.size)

# modello C_D = b + a X con X = C_L^2: equazioni normali
X, n = CL**2, CL.size
A = np.array([[np.sum(X**2), np.sum(X)], [np.sum(X), n]])
r = np.array([np.sum(X * CD), np.sum(CD)])
a, b = np.linalg.solve(A, r)
H = 4 * (n * np.sum(X**2) - np.sum(X)**2)   # hessiano di S(a, b)
print(f"k = a = {a:.5f}   C_D0 = b = {b:.5f}   H = {H:.2f} > 0")
print(f"fattore di Oswald (AR = 8): e = {1 / (np.pi * 8 * a):.3f}")
print("controllo con numpy.polyfit:",
      np.round(np.polyfit(X, CD, 1), 5))
S = np.sum((CD - b - a * X)**2)
print(f"somma dei quadrati dei residui S = {S:.3e}")
\end{python}
\begin{risultato}
k = a = 0.05000   C_D0 = b = 0.02478   H = 146.69 > 0
fattore di Oswald (AR = 8): e = 0.796
controllo con numpy.polyfit: [0.05    0.02478]
somma dei quadrati dei residui S = 2.858e-06
\end{risultato}

I dati erano stati simulati con $C_{D0}=0{,}025$ e $k=0{,}0497$ (cioè $e\approx0{,}8$ con
$AR=8$), aggiungendo un rumore casuale di ampiezza tipica $0{,}0008$: il metodo recupera i
parametri con ottima precisione (\cref{fig:minimi-quadrati}). La funzione
\texttt{polyfit} di \texttt{numpy} risolve esattamente lo stesso problema.

\begin{figure}[htb]
\centering
\includegraphics[width=\textwidth]{minimi_quadrati}
\caption{A sinistra, le misure simulate (rame) e la polare parabolica stimata (blu). A
destra, gli stessi dati nel piano $(C_{L}^{2},C_{D})$, dove il modello è una retta. I
residui, cioè i segmenti verticali grigi tra i punti e la retta, la somma dei cui quadrati
è resa minima, sono qui così piccoli (dell'ordine di $10^{-3}$) da restare quasi nascosti
dai punti (\texttt{python/fig\_minimi\_quadrati.py}).}
\label{fig:minimi-quadrati}
\end{figure}

\begin{notastorica}[Gauss, Legendre e l'asteroide perduto]
Il 1º gennaio 1801 Giuseppe Piazzi, dall'osservatorio di Palermo, scoprì Cerere, il primo
asteroide; lo seguì per poche settimane, poi lo perse nel bagliore del Sole. Il
ventiquattrenne Carl Friedrich Gauss, con un metodo di sua invenzione per determinare
l'orbita da poche osservazioni affette da errori, predisse dove ritrovarlo, e alla fine
del 1801 Cerere fu riosservato a meno di mezzo grado dalla posizione prevista. Il metodo
dei minimi quadrati fu pubblicato per la prima volta da Adrien-Marie Legendre nel 1805;
Gauss lo presentò nel 1809, affermando di usarlo dal 1795, e ne diede una giustificazione
probabilistica legata alla curva che porta il suo nome. Le condizioni del secondo ordine
per gli estremi delle funzioni di più variabili risalgono a Lagrange (1759); il
determinante delle derivate seconde porta il nome di Otto Hesse, che lo studiò nel 1844.
\end{notastorica}
